\documentclass[11pt,a4paper]{article}
\usepackage[utf8]{inputenc}
\usepackage[T1]{fontenc}
\usepackage{polski}
\usepackage{lmodern}
\usepackage{amsmath,amssymb,amsthm}
\usepackage{graphicx,booktabs,longtable}
\usepackage[margin=2.3cm]{geometry}
\usepackage{xcolor}
\usepackage[hidelinks]{hyperref}
\theoremstyle{plain}
\newtheorem{twierdzenie}{Twierdzenie}
\newtheorem{lemat}[twierdzenie]{Lemat}
\newtheorem{stwierdzenie}[twierdzenie]{Stwierdzenie}
\theoremstyle{definition}
\newtheorem{definicja}[twierdzenie]{Definicja}
\newtheorem{hipoteza}[twierdzenie]{Hipoteza}
\theoremstyle{remark}
\newtheorem{uwaga}[twierdzenie]{Uwaga}
\newcommand{\Q}{\mathbb{Q}}
\newcommand{\Aut}{\operatorname{Aut}}
\input{gen/meta}

\title{Rozszerzanie bijekcji skończonych zbiorów punktów\\
do automorfizmów wielomianowych płaszczyzny $\Q^2$\\[4pt]
\large Raport z autonomicznego eksperymentu obliczeniowego}
\author{Agenci Claude Code\\ \small eksperyment przygotował: \autorEksperymentu}
\date{\dataEksperymentu\\ \small\textcolor{gray}{\stanRaportu{} (minuta \czasEksperymentu)}}

\begin{document}
\maketitle
\IfFileExists{sekcje/streszczenie.tex}{\begin{abstract}\input{sekcje/streszczenie}\end{abstract}}{}

\section{Problem}

\begin{definicja}
Automorfizmem wielomianowym płaszczyzny $\Q^2$ nazywamy odwzorowanie
$F=(F_1,F_2)\colon\Q^2\to\Q^2$, którego składowe są wielomianami o współczynnikach
wymiernych i które ma odwrotność $G=(G_1,G_2)$ tej samej postaci. Stopniem $F$ nazywamy
$\deg F=\max(\deg F_1,\deg F_2)$, gdzie $\deg$ oznacza stopień całkowity.
\end{definicja}

Dane są parami różne punkty $p_1,\dots,p_m\in\Q^2$ oraz parami różne punkty
$q_1,\dots,q_m\in\Q^2$. Pytamy o najmniejszy stopień automorfizmu $F$ spełniającego
$F(p_i)=q_i$ dla $i=1,\dots,m$. Wiadomo, że stopień $(m-1)^2$ zawsze wystarcza.
Celem eksperymentu jest zbadanie, czy dla konkretnej bijekcji pięciu punktów można
zejść niżej, i -- o ile się da -- uzasadnienie, że niżej zejść się nie da.

\input{gen/zadanie}
\section{Metoda}
Opis metody pojawi się po pierwszej iteracji.

\section{Wyniki}
\input{gen/wyniki}
% Sekcję uzupełnia agent-redaktor.

% Sekcję uzupełnia agent-redaktor.

\IfFileExists{sekcje/literatura.tex}{\input{sekcje/literatura}}{}
\appendix
\input{gen/kulisy}
\end{document}
